\renewcommand{\chaptertitlename}{Chapter}
\chapter{Star: Enough}
\renewcommand{\chaptertitlename}{Capítulo}

\nowplaying{GLITCH IN YOUR HEART}{gabemtnz}

\lettrine{L}{a} muralla aparece de golpe, como si el mundo la hubiera escondido hasta el último segundo solo para asustarme más. Es enorme. Es antigua —piedra agrietada, jeroglíficos que nadie ha leído en siglos— y sin embargo respira: paneles cian laten entre las grietas, cables trepan como enredaderas de otra época, símbolos que no son de ningún idioma que conozca parpadean detrás del musgo. Alguien construyó esto hace mucho, muchísimo tiempo. Alguien más lo conectó después. Nadie me explicó cuándo.

\illustration{images/muro.jpg}{León y Jhonny, diminutos frente a la inmensidad de la muralla.}

León le gruñe a la muralla —\textbf{¡GRRROWWWL!}— y salta al primer saliente sin esperar respuesta, casi una orden. No hace falta que hable para que yo entienda. A nuestro alrededor, decenas ya están trepando.

Uno dibuja símbolos en el aire con un dedo y el símbolo se convierte en escalón. Otra murmura algo en un idioma antiguo y la piedra bajo sus manos se derrite y vuelve a endurecer con la forma exacta que sus dedos necesitan —alquimia, grita alguien, esto es alquimia—. Un hombre golpea la pared con dos bastones brillantes y cada golpe abre un hueco justo del tamaño de su bota, como si la muralla le debiera un favor.

Yo solo tengo mis manos.

Clavo los dedos en una grieta. Subo. Resbalo. Subo otra vez. Cada músculo que uso me informa, con bastante mala educación, que nunca antes lo había usado así. León trepa a mi lado sin esfuerzo, vuelto algo con más garras y menos vergüenza que yo, y cuando me ve resbalar por tercera vez estira una zarpa dorada y me sostiene del cuello de la sudadera como quien rescata a un gatito.

—Gracias —jadeo.

León me responde con un \textbf{¡MRRRPH!} grave y casi ofendido, como si agradecerle fuera innecesario. Por medio segundo su cara cambia: los ojos enormes, la boca abierta en una mueca ridícula, puro pánico convertido en chiste. Tan rápido como llegó, se va, y vuelve a ser mi León de siempre.

Arriba, muy arriba, hay una fila. Alguien —una figura pequeña detrás de un escritorio imposible, encajado en un saliente de la muralla como si siempre hubiera estado ahí— va anotando nombres en una pantalla que tose luz cada tres segundos. El sistema es viejísimo, lentísimo, y aun así todos hacen fila para llegar a él, para que sus nombres queden escritos en alguna parte. Como si un nombre en una pantalla fuera prueba de que uno existió.

Quiero eso. Quiero que quede constancia de que llegué hasta aquí arrastrándome.

Trepo más rápido. Los brazos ya no son míos, son un problema que cargo. Y entonces los veo.

A la derecha, un chico de pelo negro disparado hacia arriba se ríe mientras salta de saliente en saliente sin apenas tocar la piedra, empujado por algo invisible que sale de él como calor de un motor; se detiene un segundo a mitad de un salto imposible solo para estirar la mano y jalar hacia arriba a alguien que estaba a punto de caer, y sigue riendo como si nada de esto fuera grave.

Más arriba, alguien envuelto en una armadura que brilla como si tuviera una estrella metida adentro golpea la piedra con el puño una y otra vez —diez, veinte, cien golpes en el tiempo que a mí me toma dar uno solo— y cada golpe deja un peldaño nuevo. No grita. Solo respira distinto, como si estuviera ardiendo por dentro y no le importara.

A mi izquierda, un chico azul y metálico dispara algo desde su brazo hacia la piedra, calcula, dispara otra vez, y sube por los agujeros que él mismo va abriendo con una precisión que no parece nerviosa ni un solo instante, como si escalar una muralla milenaria fuera un martes cualquiera.

Y arriba de todos, sin tocar nada, flotando con las manos en los bolsillos y una expresión de aburrimiento total, un chico común, de cara redonda y ropa de escuela, sube directo hasta la fila del registro sin haber sudado ni una gota. Nadie lo aplaude. Él tampoco parece querer que lo hagan.

Todos tienen algo. Yo tengo las manos raspadas y un León.

``¿Qué soy sin nada de esto?'', pienso, y la pregunta duele más que los brazos. Ni una habilidad, ni un cuerpo entrenado, ni —hasta donde sé— una sola moneda de lo que sea que use este mundo para comprar algo. ¿Existe el dinero aquí? No lo sé. No tengo ni eso para no saberlo.

Llego a la fila. Tres personas por delante de mí. Puedo ver la pantalla, puedo ver la boca del que registra moviéndose despacio, preguntando nombres uno por uno como si el fin del mundo no estuviera cayendo detrás de nosotros en pedazos con forma de estrella.

Dos personas.

La pantalla parpadea.

Una persona.

Se apaga. No hay drama, no hay explosión: simplemente deja de estar encendida, como una vela a la que se le acabó el aire. El que registraba se queda mirando el aparato con las manos abiertas, sin nada que ofrecer. Una barrera de luz —fría, indiferente— se cierra sobre el hueco. El paso está bloqueado. Justo para mí. Justo a tres nombres de distancia del mío.

\illustration{images/Chapter-3/muralla-bloqueada.png}{Jhonny, a tres personas del registro, ve apagarse la pantalla mientras las estrellas siguen cayendo.}

Me doy vuelta buscando una salida, una escalera lateral, algo, y entonces los veo: los meteoritos, más cerca que nunca, cayendo con una lentitud que no debería ser posible a esa velocidad, iluminando la muralla entera de un naranja que no combina con el cian de los paneles. Puedo distinguir, ahora, la forma exacta de cada uno. Estrellas. Todas estrellas.

León se sienta a mi lado, en el borde de un saliente, con las patas colgando sobre el vacío como si esto fuera una tarde cualquiera. No dice nada. No hace falta.

Soy inservible. Lo pienso con esa claridad horrible que solo aparece cuando ya no hay tiempo para mentirme: a este paso no voy a tener nada. Ni una habilidad. Ni un cuerpo que responda. Ni siquiera la certeza de si el dinero existe en este lugar y, de existir, si alguna vez tendré alguno. Otros suben con fuego, con piedra líquida, con precisión de máquina, con una calma que ni siquiera necesita esforzarse. Yo subo con las uñas rotas y una pregunta que no sé responder.

Y sin embargo apoyo la mano en la melena de León —tibia, enorme, ridículamente suave para algo tan filoso— y algo en el pecho se acomoda un poco. No tengo poder. No tengo fuerza. No tengo, hasta donde puedo comprobar, absolutamente nada que valga la pena anotar en esa pantalla apagada.

Tengo esto. Tengo a alguien que baja por mí cuando resbalo, que se sienta a mi lado cuando el mundo se cierra en la cara. Quizás, solo por ahora, con eso baste.

Quizás tú también estés midiendo tu valor por lo que puedes hacer, por lo que puedes mostrar, por lo que cabe en una pantalla. Pregúntate, aunque sea un segundo, aunque duela: si hoy no pudieras hacer nada de lo que te hace sentir útil, ¿quién quedaría sentado a tu lado en el borde del abismo? Esa persona —esa compañía, ese amor que no pide currículum— también eres tú. También cuenta. También basta.

Detrás de nosotros, el cielo entero sigue cayendo.
